\section{Transformers and Multi-Head Attention}
\textit{Explain the architecture of a Transformer model and implement a minimal
Transformer encoder using PyTorch. Highlight the role of multi-head attention and
position-wise feed-forward networks. Provide code snippets to support your
answer.}

The Transformer architecture was introduced to address the
limitations of recurrent models in sequence modeling tasks,
most notably their sequential computation and difficulty
in capturing long-range dependencies.
Unlike RNN-based models, Transformers rely entirely on
attention mechanisms and dispense with recurrence and
convolution altogether~\cite{vaswaniAttentionAllYou2017}.
This allows for full parallelization during training and
more efficient modeling of global context.

A Transformer is composed of stacked encoder and decoder
blocks.
In this section, we focus on a minimal Transformer encoder,
which is sufficient to illustrate the core architectural
components.
Each encoder block consists of two main sub-layers:
a multi-head self-attention mechanism and a position-wise
feed-forward network.
Both sub-layers are wrapped with residual connections
followed by layer normalization.

Self-attention enables the model to compute contextualized
representations of tokens by attending to all other tokens
in the input sequence.
Given an input sequence represented as a matrix
$X \in \mathbb{R}^{T \times d_{\text{model}}}$, self-attention
projects the input into queries, keys, and values using
learned linear transformations
\begin{equation}
Q = XW_Q, \quad K = XW_K, \quad V = XW_V.
\end{equation}
The attention output is then computed using scaled dot-product
attention
\begin{equation}
\text{Attention}(Q, K, V) =
\text{softmax}\left(\frac{QK^\top}{\sqrt{d_k}}\right)V,
\end{equation}
where $d_k$ is the dimensionality of the key vectors.
The scaling factor stabilizes gradients for large values
of $d_k$.

Multi-head attention extends this idea by running several
attention operations in parallel.
Each head attends to the input sequence using different
learned projections, allowing the model to capture
multiple types of relationships simultaneously.
Formally, the output of multi-head attention is given by
\begin{equation}
\text{MultiHead}(Q, K, V) =
\text{Concat}(\text{head}_1, \dots, \text{head}_h)W_O,
\end{equation}
where each head is an independent attention computation.
This mechanism significantly increases the expressive
power of the model without a large computational overhead.

Following the attention sub-layer, each encoder block
contains a position-wise feed-forward network.
This network is applied independently and identically
to each position in the sequence, and consists of two
linear transformations with a non-linear activation
function in between
\begin{equation}
\text{FFN}(x) = \max(0, xW_1 + b_1)W_2 + b_2.
\end{equation}
The feed-forward network increases the representational
capacity of the model and introduces non-linearity
after the attention operation.

Since the Transformer contains no recurrence or convolution,
positional information must be injected explicitly.
This is achieved using positional encodings, which are added
to the input embeddings to provide the model with information
about token order.
A minimal Transformer encoder implemented in PyTorch is
shown in \Cref{lst:transformer}.
\begin{listing}
    \begin{minted}{Python}
class TransformerEncoderLayer(nn.Module):
    def __init__(self, d_model, num_heads, d_ff, dropout=0.1):
        super().__init__()
        self.self_attn = nn.MultiheadAttention(d_model, num_heads, dropout=dropout)
        self.ffn = nn.Sequential(
            nn.Linear(d_model, d_ff),
            nn.ReLU(),
            nn.Linear(d_ff, d_model)
        )
        self.norm1 = nn.LayerNorm(d_model)
        self.norm2 = nn.LayerNorm(d_model)
        self.dropout = nn.Dropout(dropout)

    def forward(self, x):
        # x: (seq_len, batch_size, d_model)
        attn_output, _ = self.self_attn(x, x, x)
        x = self.norm1(x + self.dropout(attn_output))
        ffn_output = self.ffn(x)
        x = self.norm2(x + self.dropout(ffn_output))
        return x
    \end{minted}
    \caption{Minimal Transformer Implmentation}\label{lst:transformer}
\end{listing}

